%% kapa: exam=1 round=24 year=2013 period=1 form=A subject=civil_law subject_ko=민법 seq=38 number=38
%% source: Q-Net 감정평가사 기출문제 (cst00309) artlSeq=5073802 file=extracted/1008/5073802/2013년 제24회 감정평가사 1차시험 기출문제/1교시(A형).hwp
%% conversion: kordoc 4.24.51 -> tools/gen1.py (auto); review=proofread-sonnet; flags=none
\begin{question}{38}
질권에 관한 설명으로 옳지 않은 것은?
\begin{choices}
  \item 지역권은 질권의 목적이 되지 못하나, 주식은 질권의 목적이 된다.
  \item 증서가 있는 채권에 대한 질권설정은 그 증서를 질권자에게 교부하여야 효력이 생긴다.
  \item 채권질권공동입질설에 따르면 책임전질은 동산질권이다.
  \item 질권자는 질물로부터 생기는 천연과실을 수취하여 다른 채권자보다 먼저 그의 채권의 변제에 충당할 수 있다.
  \item 질권의 행사는 피담보채권의 시효진행을 방해하지 않는다.
\end{choices}
\end{question}
\answer{3}
